\section{Methodology and System Design}\label{sec:methodology}

\subsection{System boundary}
We evaluate \Argorix as a systems artifact for compiled governance, not as a
security certification.  The implemented path consists of language constructs,
semantic validation, typed IR/bytecode, a fail-closed VM, an executable-provider
registry, trace generation, a SecurityReport, and an EvidenceBundle verifier.
The declarative path consists of Agent Passport fields, \ATrust and DCP-AI
metadata, external-provider contracts, bridge declarations, and local
\code{ans_name} metadata.  Proposed paths include source-digest binding,
signed bundles, append-only logs, DID/VC resolution, and Sovereign DNS/ANS
federation.

\begin{table}[t]
\centering
\caption{Representative construct families and their boundary in the studied
implementation.}
\label{tab:language-constructs}
\begin{tabularx}{\linewidth}{@{}XXX@{}}
\toprule
Construct family & Representation & Boundary \\
\midrule
Agents and messages & Typed declarations & Implemented \\
Provider contracts & Declarative metadata & Declarative \\
Policies and passports & Semantic plus runtime controls & Implemented \\
ATrust evidence maps & Validated evidence links & Declarative \\
\bottomrule
\end{tabularx}

\end{table}

\Cref{tab:language-constructs} is a method table, not a marketing claim:
represented constructs are classified by what the current artifact actually
does with them.

\subsection{Policy decision lattice}
The v0.1 snapshot exposes a design defect: all detailed policy findings are
reported as ``unknown policy rule.''  We therefore do not treat those findings
as ordinary deny/pass observations.  Instead, we specify a v0.2 typed policy
decision lattice that future artifacts must implement before policy outcomes
can be analyzed as controlled results.

\begin{table}[t]
\centering
\caption{Typed policy decision lattice required for v0.2. The current v0.1
snapshot motivates the \code{UNKNOWN_RULE} outcome; it does not yet implement
the full lattice.}
\label{tab:policy-lattice}
\begin{tabularx}{\linewidth}{@{}XXX@{}}
\toprule
Outcome & Meaning & Current status \\
\midrule
PASS & Known rule satisfied & Implemented in controlled matrix \\
DENY & Known rule violated and execution must stop & Implemented in controlled matrix \\
REVIEW & Known rule requires human review & Implemented in controlled matrix \\
UNKNOWN\_RULE & Unsupported policy identifier or configuration error & Implemented as diagnostic outcome \\
ERROR & Malformed policy object or parse/semantic failure & Implemented in controlled matrix \\
\bottomrule
\end{tabularx}

\end{table}

\Cref{tab:policy-lattice} separates known pass/deny/review outcomes from
configuration errors and malformed policy objects.  Required known rules are
\code{provider_executable_allowlist},
\code{external_provider_without_adapter}, \code{runtime_network_denied},
\code{secrets_denied}, \code{key_material_denied},
\code{passport_residency_declared}, \code{passport_residency_conflict},
\code{source_digest_present}, \code{bundle_digest_match}, and
\code{evidence_report_consistency}.  Until those outcomes are generated, the
paper reports policy evidence as a negative diagnostic finding.

\subsection{Snapshot inventory and normalization}
The observational input is an authorized snapshot of 33 request directories
from the chatbot runtime's generated output.  Publication authorization does
not relax minimization: identifiers are aggregated, excerpts are bounded, and
prompt content is not reconstructed.  A complete directory must have all five
expected artifacts listed in \cref{sec:evidence}; 27 meet that criterion and
six contain source only.  We retain incomplete directories in the denominator
and represent unavailable values as missing.

\begin{table}[t]
\centering
\caption{Committed normalized datasets. Values are regenerated from the
authorized artifact snapshot.}
\label{tab:dataset}
\begin{tabularx}{\linewidth}{@{}XXX@{}}
\toprule
Normalized source & Rows / records & Role \\
\midrule
runtime\_summary.json & 33 & Session-level normalized summary \\
sessions.csv & 33 & Tabular session observations \\
event\_counts.csv & 1890 & Per-session event-kind counts \\
verification-results.json & 27 & Verification result groups \\
\bottomrule
\end{tabularx}

\end{table}

\Cref{tab:dataset} separates session-level normalization, event-level counts,
and verifier outcomes so that totals can be recomputed without treating one
row type as another.  The analysis parses available JSON, records parse
errors, counts sizes and artifact presence, extracts execution and policy
fields, counts ledger-event kinds, and records passport, runtime-profile,
adapter, provider-boundary, and digest metadata.  A separate PowerShell
procedure invokes the repository's offline evidence verifier for every
complete bundle and stores exit code and bounded standard output.

\subsection{Controlled artifact matrix}
The snapshot alone is not a behavioral corpus.  We therefore add a controlled
artifact matrix for v0.2.  Country and residency fields are evaluated as local
metadata.  The controlled matrix covers 31 national country codes plus EU as an
optional regional residency zone.  These cases test parser validation, policy
typing, residency-conflict classification, and evidence generation.  They do
not establish legal compliance or physical storage location.

The 31-country matrix evaluates syntactic and semantic handling of
jurisdictional metadata. It does not evaluate legal compliance, physical
residency, or external identity verification.  Cases include positive controls,
negative controls, review controls, tamper controls, and unknown/error
controls.  Each case reports expected outcome, observed outcome, policy
decision, provider decision, evidence status, source digest status, trace event
count, and whether side-effect execution occurred.

\begin{table}[t]
\centering
\caption{Controlled evaluation matrix generated for v0.2. Rows describe local
metadata, policy, provider, evidence, and source-digest cases, not legal or
external identity assurance.}
\label{tab:controlled-matrix}
\begin{tabularx}{\linewidth}{@{}XXX@{}}
\toprule
Case & Expected outcome & Claim status \\
\midrule
valid simulated provider plus known policy & PASS; no external side effect & Required experiment \\
valid passport with residency declared & PASS metadata validation & Required experiment \\
valid EvidenceBundle & offline verifier passes & Observed for current bundles \\
valid source digest match & source verification passes & To be generated in v0.2 \\
external provider without executable adapter & DENY & Observed as blocked boundary \\
network attempt under denied profile & DENY & Observed as denied event \\
plaintext secret or key material attempt & DENY & Observed as denied class \\
residency mismatch or unknown jurisdiction & REVIEW & Required experiment \\
modified trace / report / bytecode & verifier fails & Required tamper experiment \\
source mismatch & source verifier fails & To be generated in v0.2 \\
unknown policy rule & UNKNOWN\_RULE configuration error & Required v0.2 diagnostic \\
malformed policy object & ERROR parse/semantic failure & Required experiment \\
\bottomrule
\end{tabularx}

\end{table}

\Cref{tab:controlled-matrix} now separates the observational runtime snapshot
from the generated v0.2 controlled matrix.  In particular, source-digest
matching and source-mismatch detection are generated by the matrix fixtures;
the historical runtime EvidenceBundles still do not verify \code{session.argx}.
The table entries are Generated in controlled matrix cases rather than inferred
from the repeated runtime snapshot.

\begin{table}[t]
\centering
\caption{Passport jurisdiction coverage in the controlled matrix. EU is not
counted as a country; it is used only as a regional residency zone.}
\label{tab:passport-coverage}
\begin{tabularx}{\linewidth}{@{}XXX@{}}
\toprule
Region group & Countries & Valid PASS cases \\
\midrule
home & 1 & 1 \\
latin\_america & 6 & 6 \\
north\_america & 2 & 2 \\
europe & 10 & 10 \\
asia\_pacific & 7 & 7 \\
africa\_middle\_east & 5 & 5 \\
\bottomrule
\end{tabularx}

\end{table}

\Cref{tab:passport-coverage} reports only metadata coverage over recognized
country codes.  It is not a statement about legal coverage, data-center
location, or external citizenship verification.

\begin{table}[t]
\centering
\caption{Controlled artifact matrix summary derived from
\code{results/controlled\_matrix.json}.}
\label{tab:matrix-summary}
\begin{tabularx}{\linewidth}{@{}XXX@{}}
\toprule
Measure & Value & Interpretation \\
\midrule
Total controlled rows & 57 & Generated deterministic fixtures \\
Countries tested & 31 & National ISO alpha-2 metadata only \\
Valid same-country passport cases & 31 & One PASS fixture per country \\
False allow / deny / review & 0 / 0 / 0 & Outcome-classification checks \\
Evidence pass / fail & 53 / 4 & Bundle or digest-consistency status \\
\bottomrule
\end{tabularx}

\end{table}

\Cref{tab:matrix-summary} is the compact paper-ready view of the generated
case table; the full row-level CSV remains the auditable source.

\subsection{Ablations and baselines}
The ablation plan follows a systems-paper discipline: remove one boundary at a
time and measure false allows, false denies, correct denies, correct reviews,
tamper detection rate, verification latency, evidence size, trace overhead,
compile time, and policy evaluation time.  These measurements are not present
in the current snapshot and must not be claimed until generated.

\begin{table}[t]
\centering
\caption{Required ablation study. Metrics in the last column are measurement
obligations, not current results.}
\label{tab:ablation}
\begin{tabularx}{\linewidth}{@{}XXX@{}}
\toprule
Variant & Removed boundary & Do not claim until measured \\
\midrule
Full ArgorixLang v0.2 & none & false allow, false deny, review accuracy, tamper rate, latency \\
w/o typed policy lattice & typed outcomes & unknown-rule collapse and aggregation errors \\
w/o executable-provider allowlist & adapter gate & external-provider false allows \\
w/o source digest & source binding & source mismatch detection \\
w/o ledger digest & trace-event ledger binding & ledger tamper detection \\
w/o report cross-check & report-to-bundle linkage & summary/detail inconsistency detection \\
w/o fail-closed default & default denial & missing-authorization false allows \\
JSON audit logs only baseline & cross-artifact verification & tamper and linkage gaps \\
\bottomrule
\end{tabularx}

\end{table}

\Cref{tab:ablation} makes the strongest future claims falsifiable: if removing
a boundary does not change a target metric, the paper should not treat that
boundary as empirically justified.

\begin{table}[t]
\centering
\caption{Baseline comparison by scope. Baselines are not adversaries; they
identify which assurance layer each system does or does not cover.}
\label{tab:baselines}
\begin{tabularx}{\linewidth}{@{}XXX@{}}
\toprule
Baseline & What it covers & Gap relative to ArgorixLang \\
\midrule
Plain application logs & chronological observations & no compiled policy or digest bundle \\
JSON audit report & structured summary & no cross-artifact verification \\
OPA/Rego-style external policy & policy decision engine & no compiled passport/evidence constructs unless integrated \\
in-toto-style provenance & signed supply-chain steps & different scope; Argorix currently unsigned local runtime evidence \\
MCP/A2A declaration only & protocol surface description & no compiled fail-closed provider boundary \\
\bottomrule
\end{tabularx}

\end{table}

\Cref{tab:baselines} keeps comparisons fair: \Argorix is scoped to language,
policy, provider boundary, and local evidence linkage, while provenance and
authorization systems solve different parts of the assurance stack.

\subsection{Analysis rules}
We define four interpretation rules in advance.  A missing artifact is
``unavailable,'' never ``passed.''  A detailed policy result overrides a
coarser summary label for policy interpretation.  Successful evidence
verification means digest and reference consistency only.  Declarative records
are not counted as externally performed actions.

We aggregate complete-session policy findings and ledger events by summing the
normalized per-session fields.  We check all four evidence digest values for
the \code{sha256:} prefix and 64 hexadecimal digits.  The verification result
is successful only when the verifier exit code is zero and its explicit
verified flag is true.

\subsection{Prompt analysis boundary}
Although the user authorized publication, we inspect only the bounded
\code{trace.injected.content} schema field.  It is absent from all 27 valid
complete traces (0 of 27 present).  The six source-only directories contain no
trace and are not assessable.  Prompt publication is a separate allowlisted
operation; its suppression is not evidence of content absence.  We therefore
perform no prompt-level qualitative analysis, quote no prompts, and report no
prompt-injection attack success rate.
